\section{Các mô hình chính xác}
\subsection{SAT với order encoding}
Cho span ứng viên $s\geq0$. Với mỗi đỉnh $v$ và $0\leq i<s$, tạo biến
\[
 x_{v,i}\Longleftrightarrow [f(v)\leq i].
\]
Dùng hằng $x_{v,-1}=\bot$ và $x_{v,s}=\top$. Các mệnh đề đơn điệu là
\begin{equation}
 \neg x_{v,i}\lor x_{v,i+1},\qquad 0\leq i<s-1.
\end{equation}
Khi đó $f(v)=a$ tương đương $x_{v,a}\land\neg x_{v,a-1}$. Với cặp đỉnh
cần cách biệt ít nhất $d$, loại mỗi cặp nhãn $(a,b)$ có $|a-b|<d$ bằng
\begin{equation}\label{eq:forbidden}
 \neg x_{u,a}\lor x_{u,a-1}\lor\neg x_{v,b}\lor x_{v,b-1}.
\end{equation}
Chọn $d=h$ cho cạnh và $d=k$ cho cặp khoảng cách hai. Các hằng tại biên được
rút gọn. Với $s=0$, một ràng buộc cách biệt dương sinh mệnh đề rỗng, đúng với
tính bất khả thi khi mọi đỉnh chỉ có thể nhận nhãn 0.

\begin{proposition}
Nếu chưa thêm ràng buộc đối xứng, CNF ở trên khả thỏa khi và chỉ khi tồn tại
một $L(h,k)$-labeling có span không quá $s$.
\end{proposition}
\begin{proof}
Một labeling trong $[0,s]$ xác định các biến ngưỡng, thỏa đơn điệu và không
chọn bất kỳ cặp nhãn bị cấm nào. Ngược lại, một dãy biến ngưỡng đơn điệu xác
định duy nhất một nhãn: chỉ số đầu tiên có giá trị đúng, hoặc $s$ nếu mọi
biến đều sai. Nếu hai nhãn vi phạm một điều kiện cách biệt, mệnh đề
\eqref{eq:forbidden} tương ứng sẽ sai. Do vậy mọi labeling giải mã đều hợp lệ.
Tịnh tiến nhãn nhỏ nhất về 0 cho phép áp dụng kết luận với định nghĩa span.
\end{proof}

Trước khi rút gọn đỉnh cố định, số biến là $|V|s$. Với $L(2,1)$, số mệnh đề
được sinh trước khi rút gọn hoặc loại trùng là
\begin{equation}
 |V|\max(s-1,0)+|E|A(s)+|D_2(G)|(s+1),
 \quad A(s)=\begin{cases}1&s=0,\\3s+1&s\geq1.\end{cases}
\end{equation}
Các số đếm trong kết quả là số biến encoder cấp phát và số mệnh đề truyền
cho solver, không phải số còn lại sau preprocessing nội bộ của SAT solver.

\subsection{ILP dạng assignment}
Cho $U$ là cận trên từ một labeling tham lam hợp lệ. Để tránh trùng ngữ nghĩa
với biến ngưỡng SAT, ký hiệu biến nhị phân ILP là $y_{v,a}$:
$y_{v,a}=1$ khi và chỉ khi $f(v)=a$. Mô hình tổng quát là
\begin{align}
 \min\quad &\ell_{\max}\\
 \sum_{a=0}^{U}y_{v,a}&=1 &&\forall v\in V,\label{eq:assignment}\\
 y_{u,a}+y_{v,b}&\leq1
 &&\forall\{u,v\}\in E,\ |a-b|<h,\\
 y_{u,a}+y_{v,b}&\leq1
 &&\forall\{u,v\}\in D_2(G),\ |a-b|<k,\\
 \sum_{a=0}^{U}a\,y_{v,a}&\leq\ell_{\max} &&\forall v\in V,\\
 y_{v,a}&\in\{0,1\},\quad \ell_{\max}\in\{0,\ldots,U\}.
\end{align}
Trong hai ràng buộc cấm, $a,b\in\{0,\ldots,U\}$. Ràng buộc
\eqref{eq:assignment} chọn đúng một nhãn. Vì mô hình không ép một đỉnh tùy ý
về 0, tối thiểu hóa $\ell_{\max}$ đạt span tối ưu nhờ bất biến tịnh tiến.
Trong triển khai $h=2,k=1$, điều kiện khoảng cách hai chỉ cần cấm nhãn bằng nhau.
Số biến là $|V|(U+1)+1$, số ràng buộc là
\[
2|V|+(3U+1)|E|+(U+1)|D_2(G)|.
\]

\begin{remark}[Quy ước nhãn bắt đầu từ 1]
Nếu dùng miền nhãn $\{1,\ldots,U'\}$ như một số cách viết assignment, tối ưu
của biến nhãn lớn nhất bằng $\lambda_{h,k}(G)+1$ trên đồ thị không rỗng.
Khi chuyển sang quy ước bắt đầu từ 0 phải dùng $U=U'-1$ và báo cáo
$\lambda=\ell_{\max}-1$. Không được đồng nhất số nhãn khả dụng với span.
\end{remark}

\subsection{ILP dạng Big-M}
Dùng biến nguyên $0\leq f_v\leq U$, $0\leq\ell_{\max}\leq U$ và
$f_v\leq\ell_{\max}$. Với mỗi cặp cần cách biệt ít nhất $d$, thêm biến
$z_{uv}\in\{0,1\}$ và
\begin{align}
 f_u-f_v&\geq d-Mz_{uv},\\
 f_v-f_u&\geq d-M(1-z_{uv}).
\end{align}
Chọn $M=U+\max(h,k)$ đủ để vô hiệu hóa bất đẳng thức không được chọn vì
$-U\leq f_u-f_v\leq U$. Mục tiêu vẫn là $\min\ell_{\max}$.
Triển khai $L(2,1)$ dùng $M=U+2$. Số biến là
$|V|+1+|E|+|D_2(G)|$; số ràng buộc là
$|V|+2|E|+2|D_2(G)|$.

Cả hai mô hình dùng nghiệm tham lam làm MIP start. Assignment loại cặp nhãn
trực tiếp nhưng kích thước tăng theo $U$; Big-M thường gọn hơn nhưng có hệ số
phụ thuộc cận trên. Đây là lý do cần thực nghiệm đối chứng, chưa phải kết luận
một mô hình luôn nhanh hơn mô hình còn lại.
